\documentclass[10pt]{article}
\usepackage[margin=1in]{geometry}
\usepackage{amsmath,amssymb}
\usepackage{booktabs}
\usepackage{array}
\usepackage{microtype}
\usepackage{xcolor}
\usepackage[hidelinks]{hyperref}

\newif\ifresultavailable
\IfFileExists{generated_results.tex}{
  \resultavailabletrue
\newcommand{\ResultCommit}{786ba9ba5c2249a710529c3b07bf82d4e29c7f36}
\newcommand{\ResultAbstract}{The locked 216-case run was classified \texttt{SUPPORTED\_WITHIN\_CONTROLLED\_CONTRACT}; the predeclared arm-specific endpoints recorded 528/528 expected outcomes. For post-result mutation, the full-versus-unbound discordance was 48 to 0 (exact two-sided McNemar $p=7.11e-15$).}
\newcommand{\ResultSection}{%
The single protected execution completed 216 cases and was
classified \texttt{SUPPORTED\_WITHIN\_CONTROLLED\_CONTRACT}. Table~\ref{tab:endpoints} reports every
predeclared arm-specific endpoint. The intervals are descriptive Wilson 95\%
intervals for the fixed matrix.

\begin{table*}[t]
\centering
\caption{Confirmatory contract endpoints. ``Expected outcome'' includes both
detection where a mechanism applies and declared non-detection where an arm has
no such mechanism.}
\label{tab:endpoints}
\begin{tabular}{@{}lcc@{}}
\toprule
Endpoint & Expected outcomes & Wilson 95\% interval \\
\midrule
Valid full acceptance & 48/48 & 0.926--1.000 \\
Valid full determinism & 48/48 & 0.926--1.000 \\
Full geometry contradiction detection & 24/24 & 0.862--1.000 \\
Full material contradiction detection & 24/24 & 0.862--1.000 \\
Full load contradiction detection & 24/24 & 0.862--1.000 \\
Full stale-hash detection & 24/24 & 0.862--1.000 \\
Hash-only semantic non-detection & 72/72 & 0.949--1.000 \\
Unbound semantic non-detection & 72/72 & 0.949--1.000 \\
Full post-mutation detection & 48/48 & 0.926--1.000 \\
Hash-only post-mutation detection & 48/48 & 0.926--1.000 \\
Unbound post-mutation non-detection & 48/48 & 0.926--1.000 \\
Full non-geometric specificity & 24/24 & 0.862--1.000 \\
Hash-only non-geometric specificity & 24/24 & 0.862--1.000 \\
\bottomrule
\end{tabular}
\end{table*}

For the paired post-result geometry mutations, the full arm alone detected
48 cases and the unbound arm alone detected
0 cases; the exact two-sided McNemar value was
$p=7.10543e-15$. This comparison describes the
presence or absence of the preregistered dependency mechanism and is not a
population-effect estimate.

On this host, median recorded successful evaluation times were
0.619 ms for full binding
($n=120$), 0.248 ms for
hash-only ($n=192$), and
0.013 ms for unbound requests
($n=216$). Failed evaluations have no elapsed-time
record, so these samples are conditional and must not be compared as balanced
performance arms. Corresponding p95 times were 1.146,
0.534, and 0.034 ms.
The median serialized full receipt was 3902 bytes and the maximum
was 4199 bytes. These are local diagnostics only.
}

}{
  \resultavailablefalse
  \newcommand{\ResultAbstract}{The confirmatory evaluation has not been executed; no result is claimed in this source state.}
  \newcommand{\ResultSection}{\textbf{Not executed.} The protected matrix remains unavailable until its preregistered resource and repository gates pass.}
  \newcommand{\ResultCommit}{not available}
}

\title{Contract-Carrying Analytical Physics Receipts for Bounded CAD Screening:\\A Controlled Conformance Study}
\author{Ryan\\Independent Research}
\date{October 2026}

\begin{document}
\maketitle

\begin{abstract}
Analytical engineering calculations are often copied away from the geometry,
material declarations, and load assumptions on which they depend. This paper
specifies and evaluates a small, fail-closed receipt that binds a versioned
analytical result to canonical CAD-primitive identity, an explicit material
card, and a named load case. The mechanism performs exact semantic checks only
where the mapping is unambiguous and can later test whether the selected
geometry has changed. We compare a full binding with a hash-only ablation and
an unbound analytical request in a preregistered synthetic conformance matrix.
\ResultAbstract{} The study concerns software-contract conformance, not physical
accuracy, solver validity, safety, or enterprise CAD/CAE generalization.
\end{abstract}

\section{Introduction}

Digital-thread and simulation-provenance research has long argued that model
artifacts, assumptions, executions, and results should remain traceable across
the engineering lifecycle \cite{hedberg2016,helu2017,nistthread,w3cprov}.
Nevertheless, a lightweight analytical calculation can still be preserved as
an isolated JSON or spreadsheet result with no executable link to the exact
geometry, material declaration, or load case that produced it. The resulting
number may remain syntactically valid after its dependencies change.

This work asks a deliberately narrow question: for closed-form calculations
over a bounded set of CAD primitives, can a machine-readable receipt detect
predeclared contradictions and post-result geometry drift while retaining
deterministic valid-case behavior? We contribute (1) a strict open reference
contract; (2) an implementation in NeuroCAD; and (3) a preregistered controlled
ablation across full, hash-only, and unbound arms. We do not claim to invent
provenance, hashing, CAD/CAE associativity, or the digital thread. We also do
not claim that correct provenance makes an analytical model physically valid.

\section{Related work and novelty boundary}

Hedberg et al. frame the digital thread around lifecycle traceability and reuse
of authoritative product data \cite{hedberg2016}. NIST reference architectures
address contextual integration among heterogeneous engineering systems
\cite{helu2017,nistthread}, while the W3C PROV data model supplies a general
representation of entities, activities, agents, and derivations
\cite{w3cprov}. Simulation-specific work captures provenance across tools and
workflow stages \cite{simprov,olaya2020,horsch2025}, and NIST has studied
traceability between systems-engineering and finite-element artifacts
\cite{nistfea}. RO-Crate provides a practical packaging model for research
objects and their metadata \cite{rocrate}.

Accordingly, the contribution here is not an algorithmic novelty claim. It is
an executable, bounded reference contract and a controlled conformance study of
one implementation choice: bind a small analytical result to canonical
primitive geometry plus explicit material and load declarations, then fail
closed on exact contradictions or later geometry drift.

\section{Receipt contract}

Let an analytical request be $R=(m,x,c)$, where $m$ is a versioned analytical
model, $x$ is a strict mapping of numeric inputs whose names encode units, and
$c$ optionally contains geometry, material, and load bindings. Canonical JSON
uses lexicographically sorted object keys, compact separators, UTF-8 encoding,
and forbids non-finite numbers. For any value $z$,

\begin{equation}
  H(z)=\operatorname{SHA256}(\operatorname{canonical\_json}(z)).
\end{equation}

The returned receipt records the executed model version, exact inputs and
outputs, disclosed assumptions, claim boundary, $H(R)$, $H(x)$, and hashes of
each declared context object. A hash establishes content identity only; it is
not a signature, an authority proof, or evidence that the content is correct.

\subsection{Geometry identity}

For one selected component, the geometry-relevant tuple is

\begin{equation}
G=(\text{id},\text{part},\text{operation},\text{shape},\text{parameters},
\text{translation},\text{rotation}).
\end{equation}

Display labels and project revision counters are excluded. Re-verification
reloads the component and returns \emph{current} exactly when its observed
$H(G)$ equals the receipt. The implementation reports closed-form isolated-
primitive volume, surface area, centroid, transformed axis-aligned bounds, and
selected unambiguous section properties. These are not Boolean part or assembly
properties. Mass is reported only for an additive primitive with an explicit
density.

\subsection{Material, load, and exact semantic checks}

A versioned material card declares a source and finite properties; a named
load case declares a coordinate system, source, and values. The evaluator does
not select materials, interpolate curves, infer directions, combine loads, or
invent safety factors. Repeated material or load values must numerically agree
with model inputs to relative tolerance $10^{-12}$.

Geometry/model checks are restricted to three exact mappings: cylinder
diameter and length for circular-shaft torsion, isolated displaced volume for
hydrostatic buoyancy, and isolated volume and total surface area for a lumped
thermal transient. Ambiguous quantities such as effective buckling length,
convection face area, and internal pipe diameter are never inferred.

\section{Experimental design}

The protocol was written before protected outcome access. The independent unit
is one generated case. Repeated evaluation of a case is a determinism check,
not an additional independent unit. The fixed seed is 20261011.

\begin{table}[t]
\centering
\caption{Preregistered confirmatory matrix.}
\label{tab:matrix}
\begin{tabular}{@{}lr@{}}
\toprule
Case family & Count \\
\midrule
Valid bound requests & 48 \\
Geometry/input contradictions & 24 \\
Material/input contradictions & 24 \\
Load/input contradictions & 24 \\
Stale expected geometry hashes & 24 \\
Post-result geometry mutations & 48 \\
Non-geometric mutations & 24 \\
\midrule
Total & 216 \\
\bottomrule
\end{tabular}
\end{table}

\begin{table}[t]
\centering
\caption{Mechanism ablation. A dash means that the arm has no such mechanism.}
\label{tab:arms}
\begin{tabular}{@{}lccc@{}}
\toprule
Capability & Full & Hash-only & Unbound \\
\midrule
Canonical component identity & yes & yes & -- \\
Semantic geometry checks & yes & -- & -- \\
Material/load agreement checks & yes & -- & -- \\
Post-result freshness API & yes & yes & -- \\
\bottomrule
\end{tabular}
\end{table}

All applicable deterministic contract endpoints require 100\% conformance; a
single miss falsifies the corresponding claim. We report exact numerators and
denominators. Wilson 95\% intervals describe the fixed matrix and are not used
as population inference. A paired exact two-sided McNemar test summarizes
full-versus-unbound post-mutation discordance, but the decision rule remains
the all-cases threshold. Timing, resident memory, and report size are host-only
diagnostics, not cross-system performance benchmarks.

The protected run requires a clean frozen Git revision, exact source and case
hashes, at least 512 MiB of free plus speculative memory, at least 2 GiB of
output storage, one process, a new output directory, and an absent execution
ledger. The ledger is written before case execution so a crash consumes the
single attempt. Runtime errors and mismatches remain outcomes; no failed case
may be removed.

\section{Results}

\ResultSection{}

\section{Threats to validity}

\paragraph{Construct validity.}
Several endpoints are intentionally contractual: a content hash should change
when hashed geometry changes. The study tests implementation conformance, not a
new physical phenomenon. An unbound arm's failure to detect drift follows from
its lack of a geometry dependency and should not be interpreted as a defect in
all unbound calculation workflows.

\paragraph{Internal validity.}
The implementer authored the code, generator, and cases. A frozen manifest,
single-attempt ledger, retained raw records, exact hashes, and fail-closed gates
limit post-outcome flexibility, but do not supply independent authorship or
adjudication.

\paragraph{External validity.}
The cases use synthetic bounded primitives and three exact geometry/model
mappings. They do not represent full CSG, assemblies, meshing, nonlinear
solvers, enterprise PLM, or heterogeneous commercial CAD/CAE environments.
No physical experiment or qualified solver comparison is included.

\paragraph{Conclusion validity.}
Wilson intervals are descriptive for this fixed matrix. Timing on one host and
small JSON artifacts cannot support a general performance claim. Any
confirmatory failure remains negative or mixed evidence under the frozen rule.

\section{Discussion}

The practical value of the contract is modest but concrete: it makes a result's
declared dependencies inspectable and gives downstream automation a fail-closed
freshness test. Its limitations are equally important. Provenance cannot repair
an inappropriate analytical model, an unreliable source, an omitted load, or
an invalid safety argument. For higher-consequence use, the receipt would need
qualified source governance, solver verification and validation, uncertainty
analysis, access control, signatures, and domain review.

\section{Reproducibility and artifact availability}

The implementation, protocol, generator, raw JSONL records, metrics, manifest,
execution ledger, and paper-generation script are retained in the NeuroCAD
repository. The result-bearing source revision is \texttt{\ResultCommit}. The
artifact records Python and platform identity, exact file hashes, case-manifest
hashes, raw-result and metric hashes, resource admission, and Git state. A
clean-environment reproduction can check software behavior, but is not labeled
independent external validation unless performed and adjudicated by an
independent party.

\section{Conclusion}

This study narrows a broad digital-thread ambition to a testable local
contract: an analytical result names and hashes the primitive geometry,
material declaration, and load case it depends on, checks exact agreements at
evaluation time, and exposes later geometry drift. The confirmatory conclusion
is limited to the preregistered synthetic matrix and the stated implementation
boundary. It does not establish physical correctness or engineering safety.

\begin{thebibliography}{9}
\bibitem{hedberg2016}
T. Hedberg Jr. et al. ``Testing the Digital Thread in Support of Model-Based
Manufacturing and Inspection.'' \emph{Journal of Computing and Information
Science in Engineering}, 16(2), 2016. \url{https://doi.org/10.1115/1.4032697}.

\bibitem{helu2017}
M. Helu, T. Hedberg Jr., and A. Feeney. ``Reference architecture to integrate
heterogeneous manufacturing systems for the digital thread.'' 2017.
\url{https://www.nist.gov/publications/reference-architecture-integrate-heterogeneous-manufacturing-systems-digital-thread}.

\bibitem{nistthread}
National Institute of Standards and Technology. ``Digital Thread for
Manufacturing.'' \url{https://www.nist.gov/programs-projects/digital-thread-manufacturing}.

\bibitem{w3cprov}
W3C. ``PROV-DM: The PROV Data Model.'' Recommendation, 2013.
\url{https://www.w3.org/TR/prov-dm/}.

\bibitem{simprov}
M. Ruscheinski et al. ``SIMPROV: Provenance capturing for simulation studies.''
\emph{PLOS ONE}, 2025. \url{https://doi.org/10.1371/journal.pone.0327607}.

\bibitem{olaya2020}
P. Olaya, J. Lofstead, and M. Taufer. ``Building Containerized Environments for
Reproducibility and Traceability of Scientific Workflows.'' 2020.
\url{https://arxiv.org/abs/2009.08495}.

\bibitem{horsch2025}
M. T. Horsch et al. ``Metadata practices for simulation workflows.''
\emph{Scientific Data}, 2025. \url{https://doi.org/10.1038/s41597-025-05126-1}.

\bibitem{nistfea}
National Institute of Standards and Technology. ``Integrating Finite Element
Analysis With Systems Engineering Models.'' 2017.
\url{https://tsapps.nist.gov/publication/get_pdf.cfm?pub_id=922742}.

\bibitem{rocrate}
RO-Crate Community. ``RO-Crate Metadata Specification 1.3.''
\url{https://www.researchobject.org/ro-crate/specification}.
\end{thebibliography}

\end{document}
